
% Page 56
% Titre de section : Alban après l'Édit de Nantes

\section*{Alban après l'Édit de Nantes}

Du calvaire de pierre, l'essentiel n'est pas d'en savoir la date de naissance, mais d'en déchiffrer les scènes qu'on a voulu y représenter. Au sommet y figure Jésus en croix. A notre gauche, c'est-à-dire à la droite du Christ, est représenté le bon larron, visage tourné vers le Divin Maître. Au-dessous, la Vierge y est accompagnée par Saint Jean et Marie Madeleine. Celle-ci est sans auréole, c'est la pécheresse. Un soldat de sa lance perce le côté de Jésus ; il en coule les gouttes de sang que recueille Nicodème dans un calice (ou graal) suivant la légende moyenâgeuse du Saint Graal. A notre droite, c'est-à-dire à la gauche du crucifié, le mauvais larron détourne son visage. Deux personnages coiffés du bonnet juif, et un troisième, tenant peut-être un pain et accosté d'une tête de bœuf, pourraient bien signifier le sacrifice ancien. Nous aurions donc, d'un côté symbolisées l'Eglise et de l'autre, la Synagogue. Un motif ornemental, peut-être dû à l'influence anglaise, est constitué par une rose aux quatre pétales en forme de cœur séparés par quatre sépales. Ce serait, prétendent les Anglais de passage, la rose des Tudor. Sur le piédroit de haut en bas, il faut lire : la descente aux enfers -- la Vierge se pâmant (=s'évanouissant) entre les bras de Saint Jean et Marie Madeleine (=N.D. du Spasme) Véronique. Sur les autres faces de la colonne, d'autres scènes : de bas en haut : l'agonie -- Jésus frappé devant Caïphe -- la Flagellation. La fuite en Egypte -- Les Rameaux -- le Jeudi Saint avec ses symboles de l'Eucharistie et du sacerdoce : agneau pascal et calice et personnage mitré. Jésus devant Pilate (=Ecce homo) -- Jésus chez Hérode -- Peut-être une Résurrection. Comment se fait-il que le sommet : côté actuellement face au mur ait été complètement abîmé ? Il comportait certainement une représentation de la Vierge entourée des quatre évangélistes ou de leur symbole. Malice du temps ? des hommes ? Les deux conjugués sans doute.
